% Source: https://artofproblemsolving.com/community/c3h606957_array_of_problems_for_warm_up_13_2009_mathcounts_handbook
% Stable identity: aops-topic-606957-q04
The small pyramid is similar to the original with volume ratio $\dfrac{1}{2}$, so the ratio of heights is $\sqrt[3]{\dfrac{1}{2}}$:
\[\dfrac{6}{h}=\dfrac{1}{\sqrt[3]2},\qquad h=\boxed{6\sqrt[3]{2}}\ \text{cm}\approx7.56\ \text{cm}.\]

\begin{asy}
usepackage("amsmath");
size(3cm);
defaultpen(fontsize(10pt));
pair A=(-3,0),B=(3,0),C=(1.5,-1.3),D=(-1.5,-1.3)+(0,0),T=(0,6); real k=2^(-1/3);
pair A2=T+k*(A-T),B2=T+k*(B-T),C2=T+k*(C-T),D2=T+k*(D-T);
draw(A--D--C--B); draw(A--B,dashed); draw(T--A); draw(T--B); draw(T--C); draw(T--D);
draw(A2--D2--C2--B2,blue); draw(A2--B2,blue+dashed); draw(T--(0,-0.65),dotted);
label("$h$",(0,1.8),E); label("$6$",T+k*((0,-0.65)-T)/2,W);
label("\mbox{}",(min(currentpicture,user=true)+max(currentpicture,user=true))/2,black); // keeps the text after the figure black when LaTeX typesets the labels inline
\end{asy}
